\begin{proposition}\label{prop:nonic117-exclusion}
An elliptic curve over a number field of degree $9$ has no point of order $117$.
\end{proposition}
\begin{proof}
Suppose that $E/K$ has such a point, where $[K:\mathbb Q]=9$.
Multiplication by $3$ gives a point of order $39$.
Let $D$ be its rational trace divisor of degree $9$ on $X=X_1(39)$, with multiplicities if the moduli field is smaller than $K$.
It is effective and supported away from the cusps.

The degree-independent $X_1(39)$ certificate, using \cite[Theorem~3.1(2), Theorem~4.13 and Table~2]{DEHMZ} and \cite[Table~1]{DvH}, gives
\[
 J_X(\mathbb Q)=[7,31122,3236688],\qquad
 J_X(\mathbb Q)[2^\infty]=[2,16],
\]
injectivity of specialization at $2$, and rational gonality $14$.
For completeness, injectivity is proved using the order-$3$ quotient
$\pi:X\to Y$ by $\langle\langle16\rangle\rangle$.
The quotient has genus $9$, finite Picard group $[6,336]\times\mathbb Z$, and a rational cuspidal specialization image whose $2$-primary order is $32$.
Since $\pi^*$ is injective on rational $2$-primary torsion, this attains the source upper bound $32$ and proves target injectivity.
The order-$3$ diamond subgroup acts trivially on $[2,16]$, whose automorphism group has order $32$.
Thus $\pi^*\pi_*=[3]$ proves source injectivity on the rational $2$-primary torsion.
Odd torsion is automatically injective at $2$.
The explicit genus-$9$ quotient, its full Picard character, the actual relation lattice and all cusp orbits are independently certified.

The global and finite source cusp signatures agree:
\[
 1^{12}2^6 12^2.
\]
Every finite cuspidal divisor of degree at most $9$ therefore lifts to a rational cuspidal divisor.
If $\overline D$ is entirely cuspidal, denote its lift by $C$.
Injectivity gives $D\sim C$, and gonality $14>9$ gives $D=C$, a contradiction.

The complete noncuspidal residue counts through degree $8$ are
\[
 (0,0,0,0,0,12,0,49).
\]
A lower-degree residue case is therefore $P+C$, where $P$ has degree $6$ or $8$ and $C$ is a globally rational cuspidal divisor of degree $3$ or $1$.
Use the point-level correspondence which transports the level point along each cyclic $5$-isogeny, and put
\[
 A_5=T_5-1-5\langle5\rangle.
\]
It annihilates $J_X(\mathbb Q)$ by Eichler--Shimura and injectivity at the odd good prime $5$.
On the cuspidal boundary, the cyclotomic Galois action $\operatorname{Fr}_5$ gives
$T_5=\operatorname{Fr}_5+5\langle5\rangle\operatorname{Fr}_5^{-1}$.
Every complete rational cusp block is fixed by this action.
Hence $A_5C=0$, also for the chosen rational base cusp, and $A_5P$ must be principal.
The previously certified tests exclude the sextic diamond orbit and four of the five octic diamond orbits.
The remaining octic orbit is the reduction of the known closed octic divisor $D_8$.
In this case injectivity and gonality give $D=D_8+C$, contradicting noncuspidal support.
Thus $\overline D$ is a single noncuspidal degree-$9$ place.
This forces the moduli field to be $K$: a trace with a nontrivial multiplicity could not have this specialization.
There is a unique prime above $2$, with ramification index $1$ and residue degree $9$.

Since this reduction is noncuspidal, the fine level-$39$ model gives good reduction of $E$ at the corresponding prime over $2$.
Its residue field is $\mathbb F_{512}$.
The point of order $117$ is prime to the residue characteristic and retains its order, so
\[
 117\mid\#E(\mathbb F_{512}).
\]
The only multiple of $117$ in the Hasse interval $[468,558]$ is $468$.

An independent enumeration of all $511$ ordinary $j$-invariants and both twists over $\mathbb F_{512}$ gives exactly three curves of order $468$.
Their $j$-invariants form one degree-$3$ Frobenius orbit.
All other curves with a point of order $39$ have order $546$ and are excluded by the preceding divisibility.
The supersingular orders are $481,513,545$, none divisible by $39$.
The ordinary census is computed by an exact Walsh transform and checked again by direct character sums for every qualified curve.

For a curve of order $468$, its rational $39$-torsion is cyclic since $\gcd(39,511)=1$.
There are $\varphi(39)/2=12$ marked level-$39$ points for each of the three eligible curves.
Thus there are exactly $36$ geometric level points over $\mathbb F_{512}$ of this type.
The known nonic curve has four distinct closed degree-$9$ diamond translates, giving all $36$ points.
Independent exact finite-field arithmetic identifies their $j$-invariants with the three eligible census entries and verifies their level points and distinct Frobenius orbits.
Consequently $\overline D=\overline {hD_0}$ for a rational diamond $h$, where $D_0$ is the known closed nonic divisor.

Specialization injectivity gives $D\sim hD_0$, and gonality gives $D=hD_0$.
In particular $D$ is closed of degree $9$, so its residue field is $K$.
The fine moduli interpretation identifies $E/K$ with the known curve, up to the corresponding number-field isomorphism.
That curve has an exact point of order $39$ and good unramified linear reduction at $29$ given by
\[
 y^2+3xy+12y=x^3+12x^2,
 \qquad \Delta=1,\qquad \#E(\mathbb F_{29})=39.
\]
The local field is $\mathbb Q_{29}$, so reduction is injective on all torsion.
Its full torsion is therefore $C_{39}$, contradicting the point of order $117$.
\end{proof}
